\subsection*{网络分区时，哪个地域可以给出首次受理}

设 g7 的权威目录和 J 由 A、B、C 三个投票副本保存，决定依第十一章的已提交前缀形成，固定配置中需要两份投票。当前承载地域 A 为第 8 代。网络分区以后，A 只能联系自己，B 与 C 可以通信。A 的旧路由缓存写着 \code{open}，B/C 所在的多数可以提交冻结和后继授予。

少数侧 A 如果仅凭本地缓存继续形成 $A_J$，多数侧又授予 B，两侧就会分别创建同一 S。协议要求首次受理提交到权威 J，A 失去多数以后，不能向客户端确认新的权威受理，也不能让外部发送者依据未提交候选付款。已经提交的原 C1 可以按平台原身份继续恢复，前提是该路径的资格与下游保护当前有效。

多数侧也要先恢复完整已提交前缀。它拥有当前写入仲裁，并不意味着其本地购买库已经物化所有旧决定。B 从 J 读到 $S1\to C1$，再核对本地投影。投影落后时返回原受理及恢复进度，扫描器补齐原数据。S2 只有在完整权威状态中未见、会话资格有效且新代次开放时，才获得新的 $A_J$。

\CFourteenFigure{ch14-original10}{分区后 B、C 多数保存原前缀并仲裁新资格，A 的旧缓存无法提交新受理，本地投影落后由 J 补齐}{fig:ch14-partition}

能够联系权威提交多数的一侧，在完成恢复和资格交接以后可以继续首次受理。另一侧可以提供明确标记新鲜度的已保存查询、引导客户端重试或等待。它没有可提交的权威决定时，新的付款入口暂时不可用。读到旧成功可以帮助恢复，读到空状态不能授权跨域首次付款。

\subsection*{把冻结、撤权和授予接成具体顺序}

目录与 J 在同一权威域内管理 g7，因此冻结命令和受理命令有明确次序。冻结位置以前已经提交的 $A_J$ 都属于旧代责任，冻结以后拒绝新的首次受理。恢复者记录 $b_{\rm cut}$，并等待目录确认其前缀完整。已经发送出去的付款单独留在待核对集合中。

随后迁移者撤销旧地域对本地购买资源的写权限。可以由资源端代次水位、受限数据库凭据与网络入口共同实施，受保护写入口逐一确认。每个水位推进等待已获锁事务结束，记录最后可恢复的物化位置。旧主失联、无法确认撤权时，若旧入口可能不经权威 J 创建新付款，后继开放条件就没有成立。

在本章 $A_J$ 先受理的模型中，旧地域失去 J 的提交能力以后，无法新增有效 C。它持有的已提交 C 都出现在 $b_{\rm cut}$ 范围内。旧进程随后即使迟到调用平台，也使用这些固定键。恢复者继续逐项核对它们，并让平台的身份保护或同域关闭覆盖这些到达。这把旧本地写入和旧外部请求分别交给对应的接收端处理。

新地域从权威目录取得第 9 代，先以 \code{recovering} 模式建立资源水位、恢复 $\Gamma$ 和分配规则、重建原身份关系，再转为 \code{restricted}。\code{restricted} 是本例的受限开放模式：旧身份只能恢复或查询，未决 C 的自动发送受各自证据条件约束，新意图则限于权威状态明确允许的范围。这个范围先要求所属分区的受理集合与会话资格完整，旧外部路径已有保护或封闭依据，再为独立的新意图登记新的有效会话。入口在同一权威受理序列中检查这份范围和新会话，才建立新的 $A_J$。全部本地投影和审计核对完成后，再进入常规 \code{open}。

\CFourteenFigure{ch14-original11}{冻结决定界定旧责任，撤销本地资格与核对外部在途分别推进，新地域先恢复证据再受限开放}{fig:ch14-recovery-order}

这里给出一次状态变迁总览。模式保存在权威域，资源端水位和物化边界有各自确认记录，恢复程序只执行满足前驱的下一步：

\begin{BookCode}
open(8,A)
  → frozen(8,b_cut)：J 禁止新的首次受理，旧责任集合固定
  → recovering(9,B)：权威分配 9，资源拒绝旧代写
      恢复完整 Γ、原身份绑定、结果、未决责任和编号规则
      对每项旧 C 确认安全继续或持久封闭的外部路径
  → restricted(9,B)：逐项恢复旧操作，接受经核验的新意图
  → open(9,B)：本地关系、覆盖边界与冲突处理全部满足准入规则
\end{BookCode}

恢复程序本身也可能崩溃。每个步骤把决定、证据和完成边界持久保存，重启后从权威状态继续。已经授予第 9 代却未物化完成时，保持 \code{recovering}。重复封闭原支付身份按同一平台身份恢复，重复物化按原 P/C/O/EV 比较和提交。新的恢复者取得更高领取代次，只改变谁有权推进这些步骤。

\subsection*{多数的交集保护了哪一份历史}

固定三副本配置中，任意两个多数至少相交一个成员。第十一章的日志匹配、提交和选举规则使这个交集把已提交受理前缀带到后继领导者~\cite{raft}。本节使用这份结果来约束目录与 J：授予第 9 代的合法后继也保存第 8 代已经提交的 $S\to C$。仅有两边各找两台机器的人数说法不足以形成这种连续性。

假如 A/B/C 全部被旧备份替换，曾经提交的后缀不在任何投票副本上，三台机器可以一致选出领导者，也可以一致查不到 C1。它们达成的是同一份旧状态。本章已经选择独立 J 跨本地回退保存，就是为了把当前恢复域之外的已受理事实继续带入后继资格。J 自身若丢失已确认历史，要先恢复它的可靠归档或其他完整证据，再执行受理开放。

复制组成员变化也要保留交接规则。把旧组三个节点直接全部换成三个空节点，然后让新组三票通过，会形成没有历史交集的仲裁域。可以采用经证明安全的成员变更协议，或在冻结下把完整已提交状态交给新组，确认恢复后再授予资格。具体复制系统负责这项变更，业务恢复者核对它交出的前缀和承载权限。

跨地域描述部署位置，安全仲裁实际依赖持久决定、参与者资格和历史传递。若目录由一个独立可靠服务提供，应用可以使用其原子条件更新来管理归属，前提是该服务的范围和持久性确实包含本次故障。若各地域只靠互相超时后自选新主，就需要另行建立能够排除双重授予的协议。

\subsection*{新代次只为新编号和新资格服务}

恢复后登记器从权威第 9 代分配 \code{E9:1}，支付受理器为新购买分配 \code{C9:1}，平台键为 \code{pay-C9:1}。旧客户端带 E1/K1 过来，服务先读原资格和绑定。E1 已关闭，旧购买入口返回 \code{operation\_expired}，授权查询沿 J 恢复 P1/C1 与原结果。E1 尚有效且 J 已有 S1，则按原 S1 恢复 C1。J 无完整覆盖时，入口报告恢复中并保留原任务。

客户端携带的地域提示（如 \code{region=B}）或路由代次参数（如 \code{epoch=9}），不能把旧 S1 变成新支付。内部 C 的编号前缀属于分配器，完整业务身份继续保持四分量。独立的新购买由用户作出新的业务决定，取得新的会话和 K，才进入新编号空间。

在 E1 尚有效、允许恢复的分支中，同 S1 来件改了金额，服务依据 J 的原输入返回 \code{conflict}。若来件保持原输入但报价现价已改为 3500，已有绑定恢复原 Q7 与 3000。E1 已关闭时，购买入口先返回过期，授权查询再恢复原购买。只有尚无绑定的新意图进入当次业务核验。
\FloatBarrier
